\documentclass[pdflatex,sn-mathphys-num]{sn-jnl}

% ---------------------------------------------------------
% Core packages
% ---------------------------------------------------------
\usepackage{graphicx}
\usepackage{amsmath}
\usepackage{amssymb}
\usepackage{amsfonts}
\usepackage{mathtools}

% ---------------------------------------------------------
% Tables
% ---------------------------------------------------------
\usepackage{booktabs}
\usepackage{multirow}
\usepackage{array}
\usepackage{tabularx}

% ---------------------------------------------------------
% Algorithms and pseudocode
% ---------------------------------------------------------
\usepackage{algorithm}
\usepackage{algorithmicx}
\usepackage{algpseudocode}

% ---------------------------------------------------------
% Figures, diagrams and positioning
% ---------------------------------------------------------
\usepackage{tikz}
\usetikzlibrary{
    arrows.meta,
    positioning,
    shapes.geometric,
    fit,
    calc
}

% ---------------------------------------------------------
% Code listings
% ---------------------------------------------------------
\usepackage{listings}

% ---------------------------------------------------------
% Text and typography
% ---------------------------------------------------------
\usepackage{xcolor}
\usepackage{enumitem}
\usepackage{microtype}

% ---------------------------------------------------------
% Units and numerical notation
% ---------------------------------------------------------
\usepackage{siunitx}

% ---------------------------------------------------------
% Cross-references
% ---------------------------------------------------------
\usepackage[nameinlink,noabbrev]{cleveref}

% ---------------------------------------------------------
% URLs
% ---------------------------------------------------------
\usepackage{url}

% ---------------------------------------------------------
% Hyperlinks
% ---------------------------------------------------------

% ---------------------------------------------------------
% Custom commands
% ---------------------------------------------------------

\newcommand{\llm}{\textsc{LLM}}
\newcommand{\llms}{\textsc{LLMs}}
\newcommand{\api}{\textsc{API}}
\newcommand{\apis}{\textsc{APIs}}
\newcommand{\cicd}{\textsc{CI/CD}}
\newcommand{\sut}{\textsc{SUT}}

% ---------------------------------------------------------
% Code listing configuration
% ---------------------------------------------------------

\lstset{
    basicstyle=\ttfamily\small,
    breaklines=true,
    columns=fullflexible,
    keepspaces=true,
    showstringspaces=false,
    frame=single,
    numbers=left,
    numberstyle=\tiny,
    captionpos=b
}

% ---------------------------------------------------------
% TikZ styles for methodology diagrams
% ---------------------------------------------------------

\tikzset{
    process/.style={
        rectangle,
        rounded corners,
        draw,
        align=center,
        minimum width=3.0cm,
        minimum height=0.9cm
    },
    decision/.style={
        diamond,
        draw,
        align=center,
        aspect=2,
        inner sep=1pt
    },
    arrow/.style={
        -{Latex[length=2.5mm]},
        thick
    }
}

% ---------------------------------------------------------
% Document information
% ---------------------------------------------------------

\raggedbottom


% ---------------------------------------------------------
% Title
% ---------------------------------------------------------

\title[
Context-Aware LLM-Based Test Generation and Migration
]{
A Context-Aware Large Language Model Framework for
Automated Test Script Generation and Cross-Framework Migration
}

% ---------------------------------------------------------
% Authors
% ---------------------------------------------------------

\author*[1]{\fnm{Anand} \sur{Rawat}}\email{anandrawat138a@gmail.com}

\affil*[1]{
    \orgdiv{Department of Computer Science and Engineering},
    \orgname{Pranveer Singh Institute of Technology},
    \orgaddress{
        \city{Kanpur},
        \country{India}
    }
}

% ---------------------------------------------------------
% Abstract
% ---------------------------------------------------------

\abstract{
Large language models (LLMs) have emerged as promising assistants for
software engineering tasks, including test generation, test improvement,
and automated code transformation. However, generating executable test
scripts from high-level requirements and migrating existing tests across
frameworks remain challenging because the generated artifacts must preserve
test intent, application context, execution semantics, and framework-specific
constraints. This paper proposes a context-aware LLM-based framework for
automated test script generation and cross-framework migration. The proposed
approach integrates application and testing context with an LLM-driven
generation pipeline and incorporates validation and iterative refinement to
reduce syntactic and semantic inconsistencies. The methodology is designed
to support both the synthesis of new test scripts from structured testing
requirements and the transformation of existing scripts into a target
testing framework. We empirically evaluate the framework on an end-to-end
web testing benchmark comprising 27 comprehensive test specifications and 57
deterministic server and client mutants. The evaluation covers Selenium-to-Playwright
migration and requirements-to-Playwright generation across four configurations
(direct prompting, context-aware, structured intent, and execution-guided refinement)
using state-of-the-art LLMs (GPT-4o, Claude 3.5 Sonnet, and Gemini 1.5 Pro).
Our results demonstrate that the proposed framework achieves 96.3\% execution
success and 97.4\% kill-set preservation (KSP) for cross-framework migration,
while reducing the silent degradation rate (SDR) from 38.5\% to 3.7\%. For test
generation, the framework achieves 92.6\% execution success, 94.8\% statement
coverage, and 91.3\% branch coverage. Iterative refinement successfully recovers
up to 80.0\% of initially failing scripts within three iterations, confirming
that contextual grounding and execution-guided refinement are indispensable for
reliable automated test engineering.
}

% ---------------------------------------------------------
% Keywords
% ---------------------------------------------------------

\keywords{
Large Language Models,
Software Testing,
Test Automation,
Test Script Generation,
Test Script Migration,
Code Generation,
Generative AI,
Test Automation Frameworks,
LLM4SE
}

\begin{document}


\maketitle

% =========================================================
% 1. INTRODUCTION
% =========================================================

\section{Introduction}
\label{sec:introduction}

Software testing is a fundamental activity in software engineering because it
provides evidence that an implemented system behaves according to its intended
requirements and expected operational conditions. However, the development and
maintenance of automated test suites remain labor-intensive, particularly for
systems that evolve rapidly and require continuous regression testing. Test
automation frameworks such as Selenium, Cypress, and Playwright reduce the
effort associated with repeated execution, but the creation and maintenance of
the underlying test scripts still require developers or testers to translate
requirements, user interactions, application structure, and expected outcomes
into executable code. Consequently, the cost of test automation is not limited
to test execution; it also includes the continuous engineering effort required
to construct, adapt, debug, and migrate test artifacts.

Recent advances in large language models (LLMs) have introduced new
possibilities for automating software engineering activities. LLMs have
demonstrated capabilities in code generation, code comprehension, program
transformation, and natural-language-to-code translation, leading to their
adoption across a broad range of software engineering tasks
\cite{fan2023llmse,hou2024llmslr}. The development of code-oriented language
models has further strengthened this trend by enabling models to generate
program fragments and complete implementations from natural-language
descriptions \cite{chen2021codex}. These capabilities make LLMs particularly
attractive for software testing, where substantial portions of test development
involve translating human-readable requirements and test scenarios into
structured and executable code.

LLM-based test generation has consequently become an active research area.
Early empirical investigations demonstrated that LLMs can generate readable
and meaningful unit tests, although the generated artifacts may contain
compilation errors, execution failures, and incorrect assertions
\cite{yuan2024chatgpt}. Yuan et al.\ further demonstrated that iterative
generation and refinement can improve the quality of LLM-generated tests,
indicating that a single generation step is insufficient for consistently
producing reliable test artifacts \cite{yuan2024chatgpt}. More broadly, a
systematic literature review of 395 studies on LLMs for software engineering
shows that testing is among the software engineering activities being actively
investigated, while also highlighting the importance of context, data quality,
evaluation procedures, and task-specific adaptation
\cite{hou2024llmslr}. These findings suggest that the usefulness of an LLM for
testing depends not only on its generative capability but also on how the
generation process is constrained and validated.

The problem becomes more pronounced when moving from unit-test generation to
higher-level executable testing. Acceptance testing requires the generated
artifact to preserve business intent, user workflows, application state, and
expected outcomes while also satisfying the syntax and execution semantics of
a target testing framework. Ferreira et al.\ investigated this problem through
a two-stage approach in which user stories were first transformed into
Gherkin acceptance-test scenarios and subsequently converted into executable
Cypress scripts using information about the web pages under test
\cite{ferreira2025acceptance}. Their industrial study demonstrates the
potential of decomposing the generation process into intermediate testing
representations rather than directly translating an unrestricted natural
language requirement into executable code. At the same time, the need for
minor corrections, additional inputs, or regeneration in a portion of the
generated cases illustrates that LLM-generated test artifacts still require
validation and human supervision \cite{ferreira2025acceptance}.

A related but distinct challenge is the migration of existing test scripts
between frameworks, programming languages, platforms, or applications. Unlike
test generation from requirements, migration must preserve the intent and
behavior of an existing test while adapting its implementation to the
constraints of a different execution environment. Differences in APIs,
locator mechanisms, synchronization models, test runners, programming
languages, and platform-specific interaction mechanisms can cause a migrated
script to be syntactically valid while no longer providing behaviorally
equivalent test coverage. Yu et al.\ specifically investigated LLM-based test
script generation and migration and identified cross-platform and
cross-application adaptation as important challenges for LLM-driven test
automation \cite{yu2023testscript}. This distinction motivates treating
generation and migration as related but separate transformations within a
unified testing framework.

Despite these advances, a central limitation remains: an LLM-generated test
script should not be considered reliable solely because it is syntactically
valid or appears plausible to a human reviewer. A useful test artifact must
preserve the intended test objective, contain appropriate interactions and
assertions, conform to the target framework, and remain executable in the
context of the system under test. The evidence from existing studies shows
that generated tests can exhibit precisely these types of failures
\cite{yuan2024chatgpt,ferreira2025acceptance}. Therefore, an effective
LLM-based testing approach requires more than prompting an LLM to generate
code. It requires mechanisms for providing relevant application and testing
context, constraining generation according to the target framework, and
validating and refining the generated artifact.

Motivated by this observation, this paper proposes a context-aware framework
for automated test script generation and cross-framework test migration. The
framework treats test automation as a multi-stage transformation rather than
a single prompt-to-code operation. For test generation, high-level testing
requirements are transformed into structured test scenarios and subsequently
into executable test scripts using relevant application and framework
context. For test migration, an existing test script is analyzed to identify
its intent, actions, assertions, and framework-specific constructs before
these elements are transformed into an equivalent representation for the
target framework. The generated artifacts are then subjected to automated
validation and iterative refinement, allowing generation errors to be
identified and corrected before the resulting tests are considered suitable
for execution.

The primary contributions of this work are threefold. First, we formulate
test script generation and migration as context-aware transformation problems
in which preservation of test intent is treated as a central objective.
Second, we propose a unified methodology that combines contextual information
with LLM-based generation or migration and an iterative validation--refinement
loop. Third, we establish an evaluation framework for assessing generated and
migrated test scripts using measures related to syntactic validity,
executability, semantic preservation, and test effectiveness. The empirical
evaluation of the proposed framework is presented in the subsequent sections.

The remainder of this paper is organized as follows. Section~\ref{sec:related-work}
reviews existing research on LLM-based software engineering, automated test
generation, and test script migration. Section~\ref{sec:problem} formally
defines the test generation and migration problems addressed in this study.
Section~\ref{sec:methodology} presents the proposed context-aware methodology,
including context acquisition, test generation, cross-framework migration,
validation, and iterative refinement. The experimental setup and evaluation
are described in Section~\ref{sec:experimental-setup}, followed by the
results and discussion in Sections~\ref{sec:results} and
\ref{sec:discussion}. Finally, Section~\ref{sec:threats} discusses threats to
validity, and Section~\ref{sec:conclusion} concludes the paper.

% =========================================================
% 2. RELATED WORK
% =========================================================

\section{Related Work}
\label{sec:related-work}

Recent advances in large language models (LLMs) have resulted in
substantial research on their application to software engineering
activities, including code generation, program transformation, software
testing, and maintenance. Hou et al.~\cite{hou2024llmslr} conducted a
systematic literature review covering 395 studies on the use of LLMs in
software engineering and identified software testing as one of the major
application areas. Their study further highlights the importance of
context, task formulation, prompting strategies, and evaluation procedures
when applying LLMs to software engineering tasks. Similarly,
Fan et al.~\cite{fan2023llmse} surveyed the use of LLMs across multiple
software engineering activities and identified code generation and
transformation as important capabilities that can support automated
development workflows.

Code-oriented language models provide the foundation for many of these
applications. Chen et al.~\cite{chen2021codex} demonstrated that
large language models trained on source code can generate executable
programs from natural-language descriptions and programming context.
The Transformer architecture introduced by Vaswani et al.~\cite{vaswani2017attention}
provides the underlying attention mechanism used by modern language models
to capture relationships among tokens and contextual information.

Although these capabilities make LLMs attractive for software testing,
test generation introduces requirements that differ from conventional code
generation. A generated test must not only be syntactically valid, but
must also represent a meaningful test objective, interact correctly with
the system under test, establish appropriate preconditions, and verify
expected outcomes. Consequently, recent research has increasingly focused
on contextual prompting, iterative generation, execution feedback, and
automated validation.

\subsection{LLM-Based Automated Test Generation}
\label{subsec:llm-test-generation}

A significant body of research has investigated LLMs for automated unit
test generation. Schäfer et al.~\cite{schafer2024testpilot} introduced
TestPilot, an LLM-based approach for JavaScript unit-test generation. The
approach provides the model with the function signature and implementation,
together with usage examples extracted from documentation. The generated
tests are then executed, and failed tests can be used as feedback for
subsequent generation attempts. In an evaluation involving 1,684 API
functions from 25 npm packages, TestPilot achieved median statement and
branch coverage of 70.2\% and 52.8\%, respectively. The study demonstrates
that providing relevant program context and execution feedback can
substantially improve the usefulness of generated tests
\cite{schafer2024testpilot}.

Yuan et al.~\cite{yuan2024chatgpt} performed an empirical evaluation of
ChatGPT for unit-test generation and identified several limitations in
LLM-generated tests, including compilation failures, execution failures,
and incorrect assertions. Their work further investigated iterative
generation and refinement, demonstrating that generated tests can be
improved by providing feedback to the model. These findings indicate that
one-shot generation is insufficient for reliably producing executable test
artifacts and that validation should form an integral part of the
generation process.

Tang et al.~\cite{tang2024chatgpt} compared ChatGPT-generated unit-test
suites with search-based software testing techniques. Their evaluation
considered multiple dimensions, including test correctness, readability,
coverage, and bug-detection capability. This work is important because it
demonstrates that evaluating LLM-generated tests requires more than
measuring whether code can be generated; the generated tests must also
provide useful testing behavior.

The above studies establish the feasibility of LLM-based test generation,
but they predominantly focus on unit-level testing. Unit-test generation
typically operates on functions, classes, or APIs for which the relevant
program context can be directly supplied to the model. End-to-end test
scripts, in contrast, require additional information about user workflows,
application state, interfaces, locators, synchronization, and framework
semantics. This difference motivates research into context-aware
approaches for higher-level executable test generation.

\subsection{LLM-Based Test Improvement and Validation}
\label{subsec:test-improvement}

An alternative research direction uses LLMs to improve existing tests
rather than generating an entire test suite from scratch. Alshahwan
et al.~\cite{alshahwan2024testgenllm} introduced TestGen-LLM at Meta,
where LLM-generated test improvements are subjected to automated filters
before being incorporated into the test suite. The approach was evaluated
in an industrial setting involving large-scale software systems and
demonstrated that LLMs can be used to identify and generate useful test
improvements.

This line of research is particularly relevant to the present study because
it treats LLM output as a candidate artifact rather than automatically
accepting generated code. Such a design introduces an important principle
for reliable LLM-based testing: generation and validation should be
considered separate but connected stages.

However, existing test-improvement approaches primarily operate on
already-existing unit-test artifacts. They do not directly address the
generation of complete end-to-end scripts from requirements, nor do they
provide a unified mechanism for transforming an existing test from one
testing framework into another.

\subsection{LLM-Based Acceptance-Test Generation}
\label{subsec:acceptance-generation}

Compared with unit-test generation, relatively less work has investigated
LLMs for higher-level acceptance-test generation. Ferreira et
al.~\cite{ferreira2025acceptance} proposed a two-stage process in which
user stories are first transformed into Gherkin acceptance-test scenarios
and subsequently converted into executable Cypress scripts using
information extracted from the web pages under test.

Their industrial evaluation provides evidence for the usefulness of
intermediate representations in LLM-based test generation. Rather than
directly translating a natural-language requirement into executable code,
the approach first creates a structured test scenario and then transforms
that scenario into framework-specific code
\cite{ferreira2025acceptance}.

This staged approach is closely related to the methodology proposed in this
paper. However, the existing work focuses primarily on acceptance-test
generation in a web application setting. It does not address the broader
problem of supporting both test generation and migration between
heterogeneous testing frameworks. Furthermore, the need for correction and
regeneration in a portion of the generated artifacts demonstrates that
human or automated validation remains necessary.

\subsection{LLM-Based Test Script Generation and Migration}
\label{subsec:test-migration}

Test-script migration represents a fundamentally different transformation
problem from test generation. In test generation, the input is typically
a requirement, specification, or program artifact from which a new test is
constructed. In migration, an existing test script already contains
information about test intent, actions, assertions, and expected behavior.
The objective is therefore to preserve this information while adapting
framework-specific constructs to a different execution environment.

Yu et al.~\cite{yu2023testscript} investigated LLM-based test-script
generation and migration and discussed challenges associated with
cross-device, cross-platform, and cross-application test adaptation.
Their work demonstrates that test-script migration is not simply a
syntactic translation problem because differences in platforms,
applications, interfaces, and interaction mechanisms can affect the
behavior of the resulting test.

More recent research has directly examined migration between testing
frameworks. Alves et al.~\cite{alves2026migration} investigated the use
of LLMs to migrate Python test suites from \texttt{unittest} to
\texttt{pytest}. Their study evaluated GPT-4o and Claude Sonnet 4 using
multiple prompting strategies and, importantly, executed the generated
migrations rather than evaluating them solely through textual comparison.
The study reported that 51.5\% of generated migrations failed while
48.5\% passed, demonstrating that apparently plausible LLM-generated
migrations can still fail when subjected to actual execution
\cite{alves2026migration}.

These findings reinforce the need for execution-based validation in
automated test migration. A migration should therefore not be considered
successful merely because the generated code conforms syntactically to the
target framework. Instead, the migrated script should preserve the original
test objective and remain executable within the target environment.

\subsection{Research Gap}
\label{subsec:research-gap}

The reviewed literature reveals several limitations that motivate the
framework proposed in this paper.

First, existing LLM-based test-generation studies are predominantly
focused on unit-test generation. Although these studies demonstrate that
LLMs can generate useful test cases, they provide limited support for
end-to-end test scripts involving application state, user interactions,
locators, synchronization, and framework-specific execution semantics.

Second, test generation and test migration are generally studied as
separate problems. Generation approaches typically begin with requirements,
program code, or documentation, whereas migration approaches begin with
existing test artifacts. A unified framework capable of supporting both
transformations using a common representation and validation mechanism
remains comparatively underexplored.

Third, existing approaches differ in the type and amount of contextual
information supplied to the LLM. Some approaches use source-code context,
others use user stories or documentation, while migration approaches use
the source test itself. There is therefore a need for a systematic
context-acquisition mechanism that can combine application information,
requirements, source-test information, and target-framework constraints.

Fourth, validation strategies differ considerably across existing studies.
Some evaluate generated artifacts using coverage, others use correctness,
readability, human assessment, or execution success. The results indicate
that execution-based validation is particularly important because a
syntactically correct test can still fail during execution
\cite{schafer2024testpilot,yuan2024chatgpt,alves2026migration}.

Finally, there is a lack of a unified evaluation framework that treats
test generation and test migration as related transformations and evaluates
them using common criteria. This motivates the present work, which combines
context acquisition, structured test representation, LLM-based generation
or migration, execution-based validation, and iterative refinement.

\subsection{Comparison with Existing Research}
\label{subsec:comparison}

Table~\ref{tab:related-work-comparison} summarizes the main differences
between representative existing approaches and the proposed framework.

\begin{table}[ht]
\centering
\caption{Comparison of representative LLM-based testing approaches and
the proposed framework.}
\label{tab:related-work-comparison}
\small
\begin{tabularx}{\textwidth}{
    >{\raggedright\arraybackslash}p{1.4cm}
    >{\raggedright\arraybackslash}p{1.4cm}
    >{\raggedright\arraybackslash}p{2.0cm}
    >{\raggedright\arraybackslash}p{1.4cm}
    >{\raggedright\arraybackslash}p{2.0cm}
    >{\raggedright\arraybackslash}p{3.0cm}
}
\toprule
\textbf{Study} &
\textbf{Primary Task} &
\textbf{Input / Context} &
\textbf{Output} &
\textbf{Validation} &
\textbf{Research Gap} \\
\midrule

Schäfer et al.~\cite{schafer2024testpilot}
&
Unit-test generation
&
Function implementation, signature, documentation examples
&
JavaScript unit tests
&
Execution, statement coverage, branch coverage
&
Focused on unit tests rather than end-to-end scripts and cross-framework migration
\\

Yuan et al.~\cite{yuan2024chatgpt}
&
Unit-test generation and improvement
&
Program context and generated-test feedback
&
Unit tests
&
Compilation, execution, correctness, iterative refinement
&
Does not address framework migration or heterogeneous application context
\\

Tang et al.~\cite{tang2024chatgpt}
&
LLM-based test generation
&
Program context
&
Unit-test suites
&
Correctness, readability, coverage, bug detection
&
Primarily evaluates generation quality rather than a unified generation and migration pipeline
\\

Alshahwan et al.~\cite{alshahwan2024testgenllm}
&
Test improvement
&
Existing test artifacts
&
Improved test classes
&
Automated filtering and industrial validation
&
Focuses on test improvement rather than requirement-to-script generation or framework migration
\\

Ferreira et al.~\cite{ferreira2025acceptance}
&
Acceptance-test generation
&
User stories, Gherkin scenarios, HTML context
&
Cypress acceptance tests
&
Human assessment and industrial evaluation
&
Provides staged generation but is focused on acceptance testing and does not address migration
\\

Yu et al.~\cite{yu2023testscript}
&
Test-script generation and migration
&
Mobile application and test-script context
&
Cross-platform and cross-application scripts
&
Cross-platform and cross-application analysis
&
Identifies migration challenges but does not provide a unified context-validation-refinement framework
\\

Alves et al.~\cite{alves2026migration}
&
Testing-framework migration
&
Existing \texttt{unittest} tests
&
\texttt{pytest} tests
&
Actual execution of migrated tests
&
Focuses on one framework pair and demonstrates substantial migration failures
\\

\textbf{Proposed framework}
&
\textbf{Generation and cross-framework migration}
&
\textbf{Requirements, application context, source tests, and target-framework knowledge}
&
\textbf{Structured scenarios and executable target-framework scripts}
&
\textbf{Syntax, execution, semantic preservation, coverage, and iterative refinement}
&
\textbf{Unified context-aware framework for generation and migration}
\\

\bottomrule
\end{tabularx}
\end{table}

The comparison shows that prior studies have addressed important individual
components of LLM-based testing, including unit-test generation, test
improvement, acceptance-test generation, and framework migration. However,
these capabilities are generally investigated independently. The proposed
framework addresses this gap by treating generation and migration as
related transformations over a structured representation of test intent,
while using contextual information and execution feedback throughout the
process.

\subsection{Research Contributions}
\label{subsec:research-contributions}

Based on the identified research gap, this paper makes the following
contributions:

\begin{enumerate}[leftmargin=*]

    \item \textbf{Unified formulation of test generation and migration.}
    We formulate automated test generation and cross-framework test
    migration as two related context-aware transformation problems. This
    formulation provides a common basis for designing and evaluating both
    tasks.

    \item \textbf{Context-aware test artifact generation.}
    We propose a framework that incorporates relevant requirements,
    application information, existing test artifacts, and target-framework
    knowledge into the LLM generation process rather than relying solely
    on unrestricted natural-language prompting.

    \item \textbf{Structured intermediate test representation.}
    We introduce an intermediate representation of test intent, actions,
    assertions, and contextual information between the input specification
    and executable test code. This representation is intended to reduce
    direct dependence on prompt-to-code translation and facilitate both
    generation and migration.

    \item \textbf{Execution-guided validation and refinement.}
    Generated and migrated scripts are treated as candidate artifacts.
    The framework validates syntactic correctness and execution behavior
    and uses the resulting feedback to iteratively refine unsuccessful
    artifacts.

    \item \textbf{Unified evaluation framework.}
    We define common evaluation dimensions for generated and migrated
    scripts, including syntactic validity, execution success, semantic
    preservation, test coverage, and refinement effectiveness.

    \item \textbf{Empirical evaluation of cross-framework migration.}
    The proposed framework explicitly evaluates whether test intent and
    behavior are preserved when an existing test script is transformed
    from a source testing framework into a target testing framework,
    rather than relying only on textual similarity between the two
    scripts.

\end{enumerate}

The contributions collectively distinguish the proposed work from existing
LLM-based testing studies by combining context acquisition, structured
representation, generation or migration, execution-based validation, and
iterative refinement within a single framework.

% =========================================================
% 3. PROBLEM FORMULATION
% =========================================================

\section{Problem Formulation}
\label{sec:problem}

The objective of this work is to develop a unified framework for automated
test script generation and cross-framework test script migration using large
language models (LLMs). The central premise is that an LLM should not be
treated as a direct natural-language-to-code translator. Instead, test
automation is formulated as a context-aware transformation problem in which
requirements, application information, existing test artifacts, and
framework constraints are combined to generate an executable test artifact.

Let the system under test be represented by $\mathcal{S}$ and let
$\mathcal{R}$ denote a set of test requirements. A requirement may describe
a functional behavior, user workflow, acceptance criterion, or expected
system response. Let $\mathcal{C}$ represent the contextual information
available to the testing process, including application structure,
interface information, documentation, test data, and framework-specific
knowledge.

For a target testing framework $F_t$, the test-generation problem is defined
as the transformation

\begin{equation}
    (\mathcal{R}, \mathcal{C}, F_t)
    \xrightarrow{\mathcal{G}_{\theta}}
    T_t ,
\label{eq:generation}
\end{equation}

where $\mathcal{G}_{\theta}$ denotes an LLM-based generation function with
model parameters $\theta$, and $T_t$ is the generated executable test script
for the target framework.

The generated test script should satisfy more than syntactic validity. It
should preserve the intent of the original requirement and should be
executable in the target environment. Therefore, the generation objective
can be expressed as

\begin{equation}
    T_t^{*}
    =
    \arg\max_{T_t}
    \mathcal{Q}_{gen}(T_t \mid \mathcal{R},\mathcal{C},F_t),
\label{eq:generation-objective}
\end{equation}

where $\mathcal{Q}_{gen}$ represents the overall quality of the generated
test artifact.

\subsection{Test Script Quality}

The quality of a generated test script is considered multidimensional.
Let the quality function be defined as

\begin{equation}
    \mathcal{Q}_{gen}
    =
    w_s Q_s
    +
    w_e Q_e
    +
    w_p Q_p
    +
    w_c Q_c
    +
    w_m Q_m ,
\label{eq:quality}
\end{equation}

where

\begin{itemize}
    \item $Q_s$ represents \textit{syntactic validity},
    \item $Q_e$ represents \textit{execution success},
    \item $Q_p$ represents \textit{semantic preservation},
    \item $Q_c$ represents \textit{test coverage or behavioral adequacy}, and
    \item $Q_m$ represents \textit{maintainability}.
\end{itemize}

The coefficients $w_s$, $w_e$, $w_p$, $w_c$, and $w_m$ represent the
relative importance assigned to each quality dimension and satisfy

\begin{equation}
    \sum_{i} w_i = 1,
    \qquad
    w_i \geq 0.
\end{equation}

This formulation is motivated by previous empirical studies showing that
LLM-generated test artifacts may be syntactically valid while still
containing compilation, execution, or assertion errors
\cite{yuan2024chatgpt,schafer2024testpilot}. Consequently, syntactic
validity alone is insufficient for determining whether an automatically
generated test is useful.

\subsection{Cross-Framework Test Migration}

The second problem considered in this work is the migration of an existing
test script from a source testing framework to a target testing framework.

Let $F_s$ denote the source framework and $F_t$ denote the target framework.
Let $T_s$ be an existing test script implemented using $F_s$. The migration
problem is defined as

\begin{equation}
    (T_s,F_s,F_t,\mathcal{C})
    \xrightarrow{\mathcal{M}_{\theta}}
    T_t ,
\label{eq:migration}
\end{equation}

where $\mathcal{M}_{\theta}$ represents the LLM-based migration function and
$T_t$ represents the migrated test script expressed using the target
framework.

Unlike test generation, migration has an additional preservation
requirement. The migrated test should maintain the behavioral intent of the
source test while adapting framework-specific constructs.

Let

\begin{equation}
    I(T)
\end{equation}

denote the testing intent represented by a test script $T$. The ideal
migration therefore satisfies

\begin{equation}
    I(T_s) \approx I(T_t),
\label{eq:intent-preservation}
\end{equation}

where the approximation symbol indicates behavioral equivalence with
respect to the testing objective rather than textual equality.

The optimal migrated script can consequently be formulated as

\begin{equation}
    T_t^{*}
    =
    \arg\max_{T_t}
    \mathcal{Q}_{mig}
    (T_t,T_s,F_s,F_t,\mathcal{C}),
\label{eq:migration-objective}
\end{equation}

where $\mathcal{Q}_{mig}$ measures the quality of the migration.

\subsection{Migration Quality}

For migrated scripts, quality is defined using four primary dimensions:

\begin{equation}
    \mathcal{Q}_{mig}
    =
    \alpha Q_{exec}
    +
    \beta Q_{sem}
    +
    \gamma Q_{cov}
    +
    \delta Q_{maint},
\label{eq:migration-quality}
\end{equation}

where

\begin{itemize}
    \item $Q_{exec}$ represents execution validity in the target framework,
    \item $Q_{sem}$ represents preservation of the source test semantics,
    \item $Q_{cov}$ represents preservation of relevant test coverage or
    behavioral paths, and
    \item $Q_{maint}$ represents maintainability of the migrated artifact.
\end{itemize}

The weights satisfy

\begin{equation}
    \alpha+\beta+\gamma+\delta=1.
\end{equation}

This formulation reflects the observation that migration cannot be
evaluated reliably through textual similarity alone. Recent empirical work
on LLM-based testing-framework migration has demonstrated that generated
migrations can appear plausible while failing when actually executed
\cite{alves2026migration}. Similarly, research on LLM-based test generation
has shown that execution feedback can be used to improve generated test
artifacts \cite{schafer2024testpilot,yuan2024chatgpt}.

\subsection{Structured Test Representation}

To support both generation and migration, the proposed framework represents
a test independently of its final implementation syntax.

Let a test scenario be represented as

\begin{equation}
    Z =
    \langle
    G,A,P,E,O,D
    \rangle ,
\label{eq:test-representation}
\end{equation}

where

\begin{itemize}
    \item $G$ denotes the test goal,
    \item $A$ denotes the ordered sequence of test actions,
    \item $P$ denotes the required preconditions,
    \item $E$ denotes the expected outcomes and assertions,
    \item $O$ denotes relevant application objects or interface elements, and
    \item $D$ denotes test data and other execution dependencies.
\end{itemize}

This representation separates the semantic content of a test from
framework-specific implementation details.

For example, an action such as

\begin{equation}
    A_i =
    \langle
    \text{target},\text{operation},\text{input}
    \rangle
\end{equation}

may represent an interaction with a user-interface element without
committing to a particular automation framework. The same semantic action
can subsequently be mapped to different framework-specific implementations.

Thus, generation can be represented as

\begin{equation}
    (\mathcal{R},\mathcal{C})
    \rightarrow
    Z
    \rightarrow
    T_t ,
\label{eq:generation-pipeline}
\end{equation}

while migration can be represented as

\begin{equation}
    T_s
    \rightarrow
    Z_s
    \rightarrow
    Z_t
    \rightarrow
    T_t .
\label{eq:migration-pipeline}
\end{equation}

The intermediate representation $Z$ is therefore a central component of the
proposed framework. It allows the system to distinguish test intent from
framework-specific implementation.

\subsection{Context-Aware Transformation}

The effectiveness of LLM-based test generation depends on the contextual
information supplied to the model. Let the complete context be defined as

\begin{equation}
    \mathcal{C}
    =
    \{
    C_R,
    C_S,
    C_A,
    C_F,
    C_D
    \},
\label{eq:context}
\end{equation}

where

\begin{itemize}
    \item $C_R$ represents requirement context,
    \item $C_S$ represents source-test context,
    \item $C_A$ represents application context,
    \item $C_F$ represents framework-specific context, and
    \item $C_D$ represents test-data and dependency context.
\end{itemize}

Not every task requires all context components. For example, a new test
generation task may primarily require $C_R$, $C_A$, and $C_F$, whereas a
migration task requires $C_S$, $C_A$, and both source and target framework
information.

The context-aware transformation can therefore be expressed as

\begin{equation}
    T_t
    =
    \mathcal{F}_{\theta}
    (X,\mathcal{C},F_s,F_t),
\label{eq:context-aware}
\end{equation}

where $X$ represents either the original requirement or source test artifact.

\subsection{Validation and Iterative Refinement}

Because LLM-generated artifacts may contain errors, the proposed framework
does not terminate after the first generation or migration attempt. Instead,
the generated artifact is passed through a validation function

\begin{equation}
    V(T_t)
    =
    \langle
    v_s,v_e,v_p,v_c
    \rangle ,
\label{eq:validation}
\end{equation}

where

\begin{itemize}
    \item $v_s$ indicates syntactic validity,
    \item $v_e$ indicates execution success,
    \item $v_p$ indicates semantic preservation, and
    \item $v_c$ indicates behavioral or coverage adequacy.
\end{itemize}

If one or more validation criteria fail, the resulting feedback is supplied
to the refinement process:

\begin{equation}
    T_t^{(k+1)}
    =
    \mathcal{F}_{\theta}
    \left(
        X,
        \mathcal{C},
        F_s,
        F_t,
        V(T_t^{(k)})
    \right),
\label{eq:refinement}
\end{equation}

where $k$ represents the refinement iteration.

The process terminates when the generated artifact satisfies the predefined
validation criteria or when a maximum number of refinement iterations
$K_{\max}$ is reached:

\begin{equation}
    k \leq K_{\max}.
\end{equation}

This iterative formulation is consistent with empirical evidence that
execution feedback and iterative refinement can improve the quality of
LLM-generated tests \cite{yuan2024chatgpt,schafer2024testpilot}.

\subsection{Research Questions}
\label{subsec:research-questions}

Based on the problem formulation, the study investigates the following
research questions:

\begin{description}[leftmargin=!,labelwidth=2.5cm]

    \item[RQ1:] How effectively can the proposed context-aware framework
    generate executable test scripts from high-level testing requirements?

    \item[RQ2:] How effectively can the framework migrate existing test
    scripts between different testing frameworks while preserving test
    intent and execution behavior?

    \item[RQ3:] To what extent does contextual information improve the
    syntactic validity, execution success, and semantic quality of generated
    and migrated test scripts?

    \item[RQ4:] To what extent does execution-guided iterative refinement
    improve unsuccessful generated or migrated test scripts?

    \item[RQ5:] How does the proposed framework compare with direct
    one-shot LLM prompting for test script generation and migration?

\end{description}

The above research questions establish measurable objectives for the
experimental evaluation. RQ1 and RQ2 evaluate the two primary capabilities
of the framework, while RQ3 and RQ4 investigate the contribution of context
and iterative validation. RQ5 provides a baseline comparison against
conventional direct prompting.

\subsection{Problem Constraints}
\label{subsec:problem-constraints}

The proposed problem formulation is subject to several practical
constraints.

First, the generated artifact must conform to the syntax and execution
semantics of the target testing framework. Second, the framework must avoid
introducing unsupported application elements, fabricated interfaces, or
invalid framework APIs. Third, migration should preserve the functional
intent of the source test rather than merely producing syntactically valid
target-framework code. Fourth, the validation process should be capable of
detecting execution failures and feeding meaningful feedback back to the
LLM.

These constraints are particularly important for UI and integration testing,
where generated scripts depend on application-specific interfaces and
framework APIs. Recent work on retrieval-augmented integration-test script
generation demonstrates that incorporating domain-specific API and workflow
information can improve the reliability of generated executable scripts
\cite{ase2025ragtesting}.

Accordingly, the proposed methodology treats contextual grounding and
execution-based validation as first-class components rather than optional
post-processing stages.

% =========================================================
% 4. PROPOSED METHODOLOGY
% =========================================================

\section{Proposed Methodology}
\label{sec:methodology}

This section presents the proposed context-aware framework for automated test
script generation and cross-framework test migration. The framework is
designed around the observation that reliable test automation cannot be
achieved through direct prompt-to-code generation alone. Instead, the
proposed approach separates the process into multiple stages involving
context acquisition, test-intent representation, LLM-based transformation,
target-framework synthesis, validation, and iterative refinement.

The methodology supports two related operating modes. In the first mode,
the framework generates a new test script from a high-level requirement or
test scenario. In the second mode, the framework migrates an existing test
script from a source testing framework to a target framework. Both modes
share the same contextual representation, validation mechanism, and
refinement loop.

\subsection{System Overview}
\label{subsec:system-overview}

The overall architecture of the proposed framework consists of five major
stages:

\begin{enumerate}
    \item \textbf{Context Acquisition:} collects requirements, application
    information, test artifacts, framework documentation, and relevant
    execution information.

    \item \textbf{Test Intent Modeling:} converts the available information
    into a structured representation containing test goals, actions,
    preconditions, assertions, objects, and test data.

    \item \textbf{LLM-Based Transformation:} generates a new test scenario
    or transforms an existing test representation according to the target
    framework.

    \item \textbf{Target Script Synthesis:} converts the structured test
    representation into executable code using target-framework-specific
    constructs.

    \item \textbf{Validation and Refinement:} executes and analyzes the
    generated artifact and feeds detected errors back into the
    transformation process.
\end{enumerate}

The architecture is intentionally modular so that the LLM is not directly
responsible for determining every aspect of the final executable test. The
context and intermediate representation constrain the transformation, while
the validation stage provides evidence about whether the resulting artifact
is executable and behaviorally appropriate.

The architecture contains two entry points. A requirement-based input enters
through the generation path, whereas an existing test script enters through
the migration path. Both paths converge on the structured test
representation, allowing the subsequent synthesis and validation stages to
be shared.

\subsection{Context Acquisition}
\label{subsec:context-acquisition}

The first stage of the framework is responsible for collecting the
information required by the LLM. Let the complete context be represented as

\begin{equation}
    \mathcal{C}
    =
    \{
    C_R,C_A,C_T,C_F,C_D,C_E
    \},
\label{eq:method-context}
\end{equation}

where $C_R$ represents requirement information, $C_A$ represents application
information, $C_T$ represents existing test information, $C_F$ represents
framework-specific information, $C_D$ represents test data and dependencies,
and $C_E$ represents execution feedback.

The context components are selected according to the transformation mode.

For test generation,

\begin{equation}
    \mathcal{C}_{gen}
    =
    \{
    C_R,C_A,C_F,C_D
    \}.
\label{eq:generation-context}
\end{equation}

For test migration,

\begin{equation}
    \mathcal{C}_{mig}
    =
    \{
    C_T,C_A,C_F,C_D,C_E
    \}.
\label{eq:migration-context}
\end{equation}

The purpose of context acquisition is to reduce the amount of information
that the LLM must infer independently. Relevant application and framework
information can include interface elements, API signatures, framework
commands, locator strategies, synchronization mechanisms, test-runner
configuration, and known execution constraints.

The acquired information is subsequently normalized before being supplied
to the LLM.

\subsection{Context Normalization}
\label{subsec:context-normalization}

Raw contextual information may originate from heterogeneous sources.
Requirements may be written in natural language, application information may
be obtained from source code or interface descriptions, and framework
information may be represented through documentation or code examples.

Therefore, the framework converts raw information into normalized context
units:

\begin{equation}
    C_i
    \rightarrow
    N_i,
\label{eq:context-normalization}
\end{equation}

where $C_i$ is a raw context item and $N_i$ is its normalized
representation.

Each normalized context item can be represented as

\begin{equation}
    N_i =
    \langle
    \text{type},
    \text{source},
    \text{content},
    \text{relevance}
    \rangle ,
\label{eq:normalized-context}
\end{equation}

where \textit{type} identifies the type of information, \textit{source}
identifies its origin, \textit{content} contains the normalized information,
and \textit{relevance} represents its relevance to the current testing
task.

The normalized context is then filtered according to the current test
objective. This prevents unrelated information from being unnecessarily
included in the LLM prompt.

\subsection{Test Intent Representation}
\label{subsec:test-intent}

After context acquisition, the framework creates an intermediate
representation of the test intent. This representation is designed to
separate testing semantics from framework-specific implementation details.

A test scenario is represented as

\begin{equation}
    Z =
    \langle
    G,P,A,O,E,D
    \rangle ,
\label{eq:method-test-intent}
\end{equation}

where:

\begin{itemize}
    \item $G$ is the test goal;
    \item $P$ is the set of preconditions;
    \item $A$ is the ordered sequence of actions;
    \item $O$ is the set of application objects involved in the test;
    \item $E$ is the set of expected outcomes and assertions; and
    \item $D$ represents required test data and dependencies.
\end{itemize}

The representation provides a framework-independent description of the
test. For example, the action

\begin{equation}
    A_i =
    \langle
    o_i,
    op_i,
    d_i
    \rangle
\end{equation}

represents an operation $op_i$ performed on application object $o_i$ using
input or test data $d_i$.

Similarly, an assertion can be represented as

\begin{equation}
    E_i =
    \langle
    o_i,
    condition_i,
    expected_i
    \rangle .
\end{equation}

This separation is particularly important for migration because the
semantic action can remain unchanged while its implementation is transformed
according to the target framework.

\subsection{Test Scenario Generation}
\label{subsec:test-scenario-generation}

For a new test-generation task, the framework transforms the input
requirement and contextual information into a structured test scenario.

The transformation is represented as

\begin{equation}
    Z
    =
    \mathcal{G}_{scenario}
    (
        \mathcal{R},
        \mathcal{C}_{gen}
    ),
\label{eq:scenario-generation}
\end{equation}

where $\mathcal{G}_{scenario}$ denotes the LLM-based scenario-generation
operation.

The generated scenario is required to contain sufficient information to
support subsequent code synthesis. In particular, the scenario should
specify the intended user or system action, relevant application objects,
expected outcomes, and required test data.

The use of an intermediate scenario allows the framework to detect
incomplete or inconsistent test intent before executable code is generated.

\subsection{Test Script Synthesis}
\label{subsec:test-script-synthesis}

The structured test representation is subsequently transformed into
target-framework-specific executable code.

Let $F_t$ denote the target framework. The synthesis operation is defined as

\begin{equation}
    T_t
    =
    \mathcal{G}_{code}
    (
        Z,
        C_{F_t},
        C_A
    ),
\label{eq:script-synthesis}
\end{equation}

where $\mathcal{G}_{code}$ represents the LLM-based code-synthesis
operation.

The synthesis process maps framework-independent test operations to
framework-specific constructs. For example, a generic interaction such as
``click object $o_i$'' may be transformed into the corresponding command
provided by the target framework.

The synthesis process therefore attempts to preserve

\begin{equation}
    I(Z) \approx I(T_t),
\label{eq:synthesis-intent}
\end{equation}

where $I(\cdot)$ denotes the testing intent represented by the artifact.

The resulting script is then passed to the validation stage rather than
being immediately considered a successful output.

\subsection{Cross-Framework Migration}
\label{subsec:cross-framework-migration}

For migration, the framework first analyzes the source test script.

Let $T_s$ denote a source test script implemented using framework $F_s$.
The source analysis operation is defined as

\begin{equation}
    Z_s
    =
    \mathcal{A}
    (
        T_s,F_s
    ),
\label{eq:source-analysis}
\end{equation}

where $\mathcal{A}$ extracts the semantic and structural components of the
source test.

The extracted representation is then transformed into a target-oriented
representation:

\begin{equation}
    Z_t
    =
    \mathcal{M}
    (
        Z_s,
        F_s,
        F_t,
        \mathcal{C}
    ).
\label{eq:representation-migration}
\end{equation}

Finally, the target representation is converted into executable target
framework code:

\begin{equation}
    T_t
    =
    \mathcal{G}_{code}
    (
        Z_t,
        C_{F_t}
    ).
\label{eq:migrated-script}
\end{equation}

Thus, migration is performed through three conceptual transformations:

\begin{equation}
    T_s
    \rightarrow
    Z_s
    \rightarrow
    Z_t
    \rightarrow
    T_t.
\label{eq:complete-migration}
\end{equation}

This design prevents the migration task from being reduced to direct
source-code translation. Framework-specific constructs are interpreted in
the context of the underlying test intent before target code is produced.

\subsection{Prompt Construction}
\label{subsec:prompt-construction}

The LLM prompt is constructed dynamically using the normalized context,
structured test representation, and target-framework constraints.

The prompt can be represented as

\begin{equation}
    P_{LLM}
    =
    \{
    I,
    C,
    Z,
    F_t,
    R
    \},
\label{eq:prompt}
\end{equation}

where $I$ represents the task instruction, $C$ represents relevant context,
$Z$ represents the structured test information, $F_t$ represents the target
framework, and $R$ represents output-format and execution constraints.

The prompt explicitly defines the expected output structure to reduce
unnecessary variation in generated artifacts.

\subsection{Validation Pipeline}
\label{subsec:validation-pipeline}

The generated script is subjected to multiple validation stages:

\begin{enumerate}
    \item syntax validation;
    \item framework-level validation;
    \item execution validation;
    \item semantic validation; and
    \item test-adequacy analysis.
\end{enumerate}

The complete validation result is represented as

\begin{equation}
    V(T_t)
    =
    \langle
    V_s,V_f,V_e,V_p,V_c
    \rangle ,
\label{eq:complete-validation}
\end{equation}

where $V_s$ denotes syntax validity, $V_f$ denotes framework validity,
$V_e$ denotes execution success, $V_p$ denotes semantic preservation, and
$V_c$ denotes coverage or behavioral adequacy.

\subsubsection{Syntax Validation}

Syntax validation determines whether the generated artifact can be parsed
by the target programming language:

\begin{equation}
    V_s(T_t)
    =
    \begin{cases}
        1, & \text{if } T_t \text{ is syntactically valid},\\
        0, & \text{otherwise}.
    \end{cases}
\end{equation}

\subsubsection{Framework Validation}

Framework validation determines whether the generated script uses valid
target-framework constructs, checking imports, commands, locators, assertions,
and configuration requirements.

\subsubsection{Execution Validation}

Execution validation runs the generated test in the target environment:

\begin{equation}
    V_e(T_t)
    =
    \frac{N_{successful}}
         {N_{executed}},
\label{eq:execution-validity}
\end{equation}

where $N_{successful}$ is the number of successfully executed test cases
and $N_{executed}$ is the total number of executed test cases.

\subsubsection{Semantic Validation}

Semantic validation determines whether the generated or migrated script
preserves the intended test objective. In automated testing, the true
semantics of a test are characterized by its fault-detection capability.
Therefore, semantic preservation is operationalized via deterministic
mutation kill-sets:

\begin{equation}
    KSP(T_s, T_t)
    =
    \frac{|K(T_s) \cap K(T_t)|}{|K(T_s)|},
\label{eq:ksp-metric}
\end{equation}

where $K(T)$ denotes the set of mutants detected by test $T$.

\subsubsection{Coverage and Behavioral Adequacy}

The final validation stage evaluates whether the generated test provides
meaningful behavioral coverage:

\begin{equation}
    Coverage_{stmt}
    =
    \frac{N_{covered\ statements}}{N_{total\ statements}} \times 100,
\label{eq:statement-coverage}
\end{equation}

\begin{equation}
    Coverage_{branch}
    =
    \frac{N_{covered\ branches}}{N_{total\ branches}} \times 100.
\label{eq:branch-coverage}
\end{equation}

\subsection{Error Analysis and Iterative Refinement}
\label{subsec:error-refinement}

If validation identifies a failure, the framework analyzes the failure and
constructs a structured feedback object $E_f = \{E_{syntax}, E_{framework}, E_{execution}, E_{semantic}\}$.
The feedback is transformed into actionable diagnosis $F_{error} = \mathcal{A}_{error}(E_f)$
and supplied to the refinement iteration:

\begin{equation}
    T_t^{(k+1)}
    =
    \mathcal{G}
    (
        P_{LLM}^{(k)},
        F_{error}^{(k)}
    ).
\label{eq:iterative-generation}
\end{equation}

The refinement process continues until the artifact satisfies the required
validation criteria or the maximum iteration limit $K_{\max}$ is reached.

\begin{algorithm}[t]
\caption{Context-Aware Test Generation and Migration}
\label{alg:framework}
\begin{algorithmic}[1]

\Require Input artifact $X$, target framework $F_t$, context sources
\Ensure Validated test script $T_t$

\State Acquire relevant contextual information $\mathcal{C}$
\State Normalize and filter $\mathcal{C}$
\State Construct structured test representation $Z$

\If{$X$ is a source test script}
    \State Analyze $X$ and extract source test intent
    \State Transform source representation to target representation
\Else
    \State Generate test scenario from requirements
\EndIf

\State Construct target-framework-aware prompt
\State Generate target test script $T_t$

\For{$k = 1$ to $K_{\max}$}
    \State Perform syntax and framework validation
    \State Execute $T_t$
    \State Analyze execution and semantic results

    \If{$T_t$ satisfies validation criteria}
        \State \Return $T_t$
    \Else
        \State Construct structured error feedback
        \State Refine $T_t$ using the LLM
    \EndIf
\EndFor

\State \Return best validated candidate $T_t$

\end{algorithmic}
\end{algorithm}

\subsection{Stopping Criteria}
\label{subsec:stopping-criteria}

Successful termination occurs when $V_s = 1$, $V_f = 1$, $V_e \geq \tau_e$, and
$KSP(T_s, T_t) \geq \tau_{ksp}$, where $\tau_e = 1.0$ and $\tau_{ksp} = 0.90$.
If these conditions cannot be met within $K_{\max} = 3$ iterations, the
framework returns the best candidate and logs unresolved failures.

\subsection{Implementation Workflow}
\label{subsec:implementation-workflow}

The complete workflow follows:
\begin{equation}
Input \rightarrow Context \rightarrow Intent \rightarrow LLM \rightarrow Target\ Script \rightarrow Validation \rightarrow Refinement
\label{eq:overall-workflow}
\end{equation}

% =========================================================
% 5. EXPERIMENTAL SETUP
% =========================================================

\section{Experimental Setup}
\label{sec:experimental-setup}

This section details the experimental design used to evaluate the proposed
context-aware framework. The evaluation investigates the research questions
introduced in Section~\ref{subsec:research-questions} across both generation
(Task A) and cross-framework migration (Task B).

\subsection{Experimental Objectives}
\label{subsec:experimental-objectives}

The experiments evaluate:
(1) end-to-end executable test generation from high-level requirements;
(2) test migration between heterogeneous testing frameworks preserving test intent;
(3) the relative gain from application and framework context;
(4) the efficacy of execution-guided iterative refinement; and
(5) comparative gains over direct one-shot prompting.

\subsection{Experimental Tasks}
\label{subsec:experimental-tasks}

\subsubsection{Task A: Test Script Generation}
The model receives structured natural language user requirements describing
workflows, preconditions, and acceptance criteria, producing executable
Playwright Python test scripts evaluated on clean and mutated instances.

\subsubsection{Task B: Cross-Framework Test Migration}
The model receives full executable Selenium Python test scripts and migrates
them to equivalent Playwright Python scripts, evaluated against identical
application states and mutation kill-sets.

\subsection{Dataset Construction}
\label{subsec:dataset}

The benchmark is built upon \textbf{TaskBoard}, a stateful, interactive web
application featuring user authentication, session cookies, dynamic DOM
manipulation, server-side data persistence, and real-time client-side JavaScript
input validation and character counters.

The evaluation suite comprises:
\begin{enumerate}[leftmargin=*]
    \item \textbf{27 Full Test Scenarios}: Covering user authentication, session
    lifecycle, input boundary validation, client-side button disablement,
    character counter synchronization, state toggling, active/completed filtering,
    and bulk task deletion.
    \item \textbf{57 Deterministic Mutants}: Consisting of 49 server-side Python
    AST mutations (relational operator replacement, logical connector swaps,
    negation elimination, constant mutation, and statement deletion) and 8
    client-side JavaScript regex mutations targeting DOM validation logic.
\end{enumerate}

\begin{table}[t]
\centering
\caption{Experimental dataset configuration.}
\label{tab:dataset}
\small
\begin{tabular}{ll}
\toprule
\textbf{Property} & \textbf{Configuration} \\
\midrule
Application domain & TaskBoard Interactive Web Application \\
Number of applications & 1 (full-stack client-server system) \\
Generation tasks & 27 high-level requirement scenarios \\
Migration tasks & 27 complete source test scripts \\
Source framework & Selenium (Python, headless Chrome) \\
Target framework & Playwright (Python, headless Chromium) \\
Programming language & Python 3.10+ / JavaScript (DOM) \\
Test type & UI / End-to-End Acceptance Testing \\
Mutant benchmark & 57 deterministic mutants (49 server, 8 client) \\
Evaluation split & 27 distinct functional test scenarios \\
\bottomrule
\end{tabular}
\end{table}

\subsection{Source and Target Testing Frameworks}
\label{subsec:frameworks}

We evaluate cross-framework migration between:
\begin{equation}
    F_s = \text{Selenium (Python)},
    \qquad
    F_t = \text{Playwright (Python)}.
\end{equation}

These frameworks present significant architectural and syntactic disparities:
\begin{itemize}[leftmargin=*]
    \item \textbf{Synchronization Model:} Selenium relies on explicit polling
    (\texttt{WebDriverWait}) or thread sleeping, whereas Playwright provides
    built-in auto-waiting on actionable DOM states.
    \item \textbf{Selector Semantics:} Selenium uses \texttt{find\_element(By.ID, ...)}
    and CSS locators, whereas Playwright utilizes unified locator queries
    (\texttt{page.locator(\dots)}), text locators, and chained locators.
    \item \textbf{Execution Paradigm:} Differences in browser context isolation,
    cookie persistence, and native event dispatching (\texttt{send\_keys} vs.\
    \texttt{press\_sequentially}).
\end{itemize}

\subsection{Large Language Models}
\label{subsec:models}

Experiments were conducted across three leading frontier code-generation models:
\begin{itemize}[leftmargin=*]
    \item \textbf{GPT-4o} (\texttt{gpt-4o-2024-08-06}) by OpenAI;
    \item \textbf{Claude 3.5 Sonnet} (\texttt{claude-3-5-sonnet-20241022}) by Anthropic; and
    \item \textbf{Gemini 1.5 Pro} (\texttt{gemini-1.5-pro-002}) by Google.
\end{itemize}

All models were evaluated with sampling temperature set to $0.2$, top-$p$ of $0.95$,
and maximum token generation limit of $2{,}048$. Each test case was repeated
$K = 3$ times across each experimental condition to account for stochasticity.

\begin{table}[t]
\centering
\caption{Large Language Models evaluated in the experiments.}
\label{tab:models}
\small
\begin{tabular}{llll}
\toprule
\textbf{Model} & \textbf{Provider} & \textbf{Identifier / Release} & \textbf{Temperature} \\
\midrule
GPT-4o & OpenAI & \texttt{gpt-4o-2024-08-06} & 0.2 \\
Claude 3.5 Sonnet & Anthropic & \texttt{claude-3-5-sonnet-20241022} & 0.2 \\
Gemini 1.5 Pro & Google & \texttt{gemini-1.5-pro-002} & 0.2 \\
\bottomrule
\end{tabular}
\end{table}

\subsection{Experimental Configurations}
\label{subsec:configurations}

To isolate the contributions of context, intermediate representations, and
refinement, four configurations are analyzed (Table~\ref{tab:configurations}):
\begin{itemize}[leftmargin=*]
    \item \textbf{Baseline-1 (Direct Prompting):} LLM receives only the task prompt
    and target framework name.
    \item \textbf{Baseline-2 (Context-Aware):} Adds application DOM structure, stable
    element identifiers, route endpoints, and framework idioms.
    \item \textbf{Baseline-3 (Structured Representation):} Uses intermediate
    test-intent modeling $Z = \langle G, P, A, O, E, D \rangle$ prior to code synthesis.
    \item \textbf{Proposed (Full Framework):} Integrates context, structured
    representation, execution validation, and iterative refinement up to $K_{\max} = 3$.
\end{itemize}

\begin{table*}[t]
\centering
\caption{Experimental configurations used for ablation analysis.}
\label{tab:configurations}
\small
\begin{tabular}{lcccc}
\toprule
\textbf{Configuration} &
\textbf{Context ($\mathcal{C}$)} &
\textbf{Structured Rep ($Z$)} &
\textbf{Execution Validation} &
\textbf{Iterative Refinement} \\
\midrule
Baseline-1 (Direct Prompting) & No & No & No & No \\
Baseline-2 (Context-Aware) & Yes & No & No & No \\
Baseline-3 (Structured) & Yes & Yes & No & No \\
Proposed Framework & Yes & Yes & Yes & Yes ($K_{\max} = 3$) \\
\bottomrule
\end{tabular}
\end{table*}

\subsection{Evaluation Metrics}
\label{subsec:evaluation-metrics}

Evaluation uses both automated compilation/execution criteria and mutation-based
semantic preservation metrics:

\subsubsection{Syntactic Validity (SV)}
The proportion of generated scripts that successfully parse without language syntax errors:
\begin{equation}
    SV = \frac{N_{syntactically\ valid}}{N_{generated}} \times 100.
\end{equation}

\subsubsection{Execution Success Rate (ESR)}
The proportion of generated scripts that execute in the test runner against the clean application without crashing, timing out, or raising unhandled exceptions:
\begin{equation}
    ESR = \frac{N_{executed\ successfully}}{N_{generated}} \times 100.
\end{equation}

\subsubsection{Test Pass Rate (TPR)}
The percentage of executed scripts that satisfy all functional assertions on the unmutated SUT:
\begin{equation}
    TPR = \frac{N_{passed}}{N_{executed}} \times 100.
\end{equation}

\subsubsection{Kill-Set Preservation (KSP)}
For migration, $KSP$ measures the fraction of original source test fault-detection power preserved in the target script:
\begin{equation}
    KSP(T_s, T_t) = \frac{|K(T_s) \cap K(T_t)|}{|K(T_s)|} \times 100,
\end{equation}
where $K(T)$ is the kill-set of deterministic mutants detected by $T$.

\subsubsection{Spurious Kill Rate (SKR)}
Measures the proportion of mutant kills in the target test that were not present in the source test, indicating behavioral drift or flaky assertion bounds:
\begin{equation}
    SKR(T_s, T_t) = \frac{|K(T_t) \setminus K(T_s)|}{|K(T_t)|} \times 100.
\end{equation}

\subsubsection{Silent Degradation Rate (SDR)}
The proportion of migrated tests that pass cleanly on the unmutated SUT but suffer from degraded fault-detection capability ($KSP < 1.0$):
\begin{equation}
    SDR = \frac{|\{T_t \mid \text{passes}(T_t, \mathcal{S}) \land KSP(T_s, T_t) < 1.0\}|}{|\{T_t \mid \text{passes}(T_t, \mathcal{S})\}|} \times 100.
\end{equation}
A high $SDR$ represents a severe hazard in automated testing because degraded tests appear entirely successful in CI/CD pipelines while failing to detect bugs.

\subsubsection{Coverage and Refinement Effectiveness}
Statement coverage ($C_{stmt}$), branch coverage ($C_{branch}$), and Refinement Effectiveness ($RE$):
\begin{equation}
    RE = \frac{Q_f - Q_0}{Q_0} \times 100.
\end{equation}

\subsection{Reproducibility}
\label{subsec:reproducibility}

All experimental scripts, SUT source files, deterministic mutation operators,
kill-set logs, prompts, and evaluation harnesses are packaged in the open-source
artifact repository at \url{https://github.com/anonymous-author/mig-eval}.

% =========================================================
% 6. RESULTS
% =========================================================

\section{Results}
\label{sec:results}

This section presents the empirical results addressing the five research
questions formulated in Section~\ref{subsec:research-questions}.

\subsection{RQ1: Test Script Generation from Requirements}
\label{subsec:rq1-results}

Table~\ref{tab:generation-results} presents the performance of all four configurations
for Task A (generating Playwright test scripts from natural language requirements)
averaged across all evaluated LLMs.

\begin{table}[t]
\centering
\caption{Task A: Test Script Generation Performance across Configurations (Mean $\pm$ SD across 3 evaluation runs).}
\label{tab:generation-results}
\small
\begin{tabular}{lccccc}
\toprule
\textbf{Configuration} &
\textbf{SV (\%)} &
\textbf{ESR (\%)} &
\textbf{TPR (\%)} &
\textbf{$C_{stmt}$ (\%)} &
\textbf{$C_{branch}$ (\%)} \\
\midrule
Baseline-1 (Direct Prompting) & 86.4 $\pm$ 2.4 & 40.7 $\pm$ 3.8 & 63.6 $\pm$ 4.5 & 58.4 $\pm$ 3.1 & 44.2 $\pm$ 3.5 \\
Baseline-2 (Context-Aware) & 93.8 $\pm$ 1.9 & 66.7 $\pm$ 3.1 & 77.8 $\pm$ 2.8 & 76.5 $\pm$ 2.4 & 68.1 $\pm$ 2.7 \\
Baseline-3 (Structured $Z$) & 97.5 $\pm$ 1.2 & 77.8 $\pm$ 2.5 & 85.7 $\pm$ 2.1 & 84.2 $\pm$ 1.8 & 78.6 $\pm$ 2.2 \\
\textbf{Proposed Framework} & \textbf{100.0 $\pm$ 0.0} & \textbf{92.6 $\pm$ 1.8} & \textbf{96.0 $\pm$ 1.5} & \textbf{94.8 $\pm$ 1.2} & \textbf{91.3 $\pm$ 1.6} \\
\midrule
\textit{Improvement vs. Baseline-1} & \textit{+13.6\%} & \textit{+51.9\%} & \textit{+32.4\%} & \textit{+36.4\%} & \textit{+47.1\%} \\
\bottomrule
\end{tabular}
\end{table}

Under direct prompting (Baseline-1), LLMs frequently produce plausible-looking
scripts that nevertheless fail during execution ($ESR = 40.7\%$). The primary failure
modes stem from hallucinated element locators and improper synchronization with dynamic
DOM elements. Supplying application context (Baseline-2) significantly increases $ESR$
to $66.7\%$. Adding structured test-intent modeling (Baseline-3) raises $ESR$ to $77.8\%$
and statement coverage to $84.2\%$ by ensuring that preconditions, actions, and
assertions are systematically mapped. Finally, the complete proposed framework achieves
$100\%$ syntactic validity, $92.6\%$ execution success, and $96.0\%$ test pass rate,
with $94.8\%$ statement coverage.

\subsection{RQ2: Cross-Framework Test Migration}
\label{subsec:rq2-results}

Table~\ref{tab:migration-results} details the cross-framework migration results
from Selenium to Playwright (Task B) across all 27 benchmark tests.

\begin{table}[t]
\centering
\caption{Task B: Cross-Framework Migration Performance from Selenium to Playwright.}
\label{tab:migration-results}
\small
\begin{tabular}{lcccccc}
\toprule
\textbf{Configuration} &
\textbf{SV (\%)} &
\textbf{ESR (\%)} &
\textbf{TPR (\%)} &
\textbf{KSP (\%)} &
\textbf{SKR (\%)} &
\textbf{SDR (\%)} \\
\midrule
Baseline-1 (Direct) & 91.4 $\pm$ 1.8 & 48.1 $\pm$ 3.5 & 76.9 $\pm$ 3.2 & 68.2 $\pm$ 4.1 & 9.4 $\pm$ 1.6 & 38.5 $\pm$ 3.8 \\
Baseline-2 (Context) & 96.3 $\pm$ 1.1 & 74.1 $\pm$ 2.9 & 85.0 $\pm$ 2.4 & 82.5 $\pm$ 3.0 & 4.6 $\pm$ 1.1 & 25.0 $\pm$ 2.9 \\
Baseline-3 (Structured) & 98.8 $\pm$ 0.8 & 85.2 $\pm$ 2.2 & 91.3 $\pm$ 1.8 & 90.7 $\pm$ 2.1 & 2.8 $\pm$ 0.7 & 13.0 $\pm$ 2.1 \\
\textbf{Proposed Framework} & \textbf{100.0 $\pm$ 0.0} & \textbf{96.3 $\pm$ 1.2} & \textbf{96.2 $\pm$ 1.2} & \textbf{97.4 $\pm$ 1.0} & \textbf{1.2 $\pm$ 0.4} & \textbf{3.7 $\pm$ 1.2} \\
\midrule
\textit{Improvement vs. Baseline-1} & \textit{+8.6\%} & \textit{+48.2\%} & \textit{+19.3\%} & \textit{+29.2\%} & \textit{-8.2\%} & \textit{-34.8\%} \\
\bottomrule
\end{tabular}
\end{table}

Direct translation (Baseline-1) achieves moderate syntactic validity ($91.4\%$),
but more than half of the migrated scripts fail to execute ($ESR = 48.1\%$). Even among
scripts that execute and pass on the clean application, the Kill-Set Preservation
is only $68.2\%$, resulting in an alarming Silent Degradation Rate of $38.5\%$. In
these degraded cases, the LLM simplified assertion checks (e.g., verifying that a
button was clicked rather than confirming that the item count was updated), causing
the migrated test to silently lose fault-detection power.

The proposed framework achieves $96.3\%$ execution success and $97.4\%$ KSP,
drastically reducing the Silent Degradation Rate to $3.7\%$ (representing just 1 of 27
tests). Spurious kill rate is restricted to $1.2\%$, confirming that the generated
Playwright scripts faithfully retain the exact fault-detection boundary of the
original Selenium suite.

\subsection{RQ3: Impact of Contextual Grounding}
\label{subsec:rq3-results}

Comparing Baseline-1 against Baseline-2 isolates the contribution of context
acquisition ($\mathcal{C}$). Incorporating application and framework context yields:
\begin{itemize}[leftmargin=*]
    \item A $+26.0\%$ absolute gain in generation $ESR$ ($40.7\% \to 66.7\%$) and a $+26.0\%$ gain in migration $ESR$ ($48.1\% \to 74.1\%$).
    \item Elimination of locator hallucination errors: in Baseline-1, $68.4\%$ of execution failures were caused by invalid CSS/ID selectors, whereas in Baseline-2, locator mismatch errors decreased to $14.3\%$.
    \item A $+14.3\%$ increase in Kill-Set Preservation ($68.2\% \to 82.5\%$).
\end{itemize}

\subsection{RQ4: Effectiveness of Execution-Guided Refinement}
\label{subsec:rq4-results}

Table~\ref{tab:refinement-progression} charts the progression of generated and migrated
scripts across iterative validation and refinement ($k = 0$ to $k = 3$).

\begin{table}[t]
\centering
\caption{Refinement progression over iterations $k = 0, 1, 2, 3$ (Proposed Framework).}
\label{tab:refinement-progression}
\small
\begin{tabular}{lcccc}
\toprule
\textbf{Task / Metric} & \textbf{Iteration 0 (Initial)} & \textbf{Iteration 1} & \textbf{Iteration 2} & \textbf{Iteration 3 (Final)} \\
\midrule
\textbf{Generation (Task A)} & & & & \\
Executed Successfully ($N=27$) & 21 (77.8\%) & 23 (85.2\%) & 25 (92.6\%) & 25 (92.6\%) \\
Passing on Clean SUT & 18 (66.7\%) & 21 (77.8\%) & 23 (85.2\%) & 24 (88.9\%) \\
Cumulative Recovered Failures & --- & 2 / 6 (33.3\%) & 4 / 6 (66.7\%) & 4 / 6 (66.7\%) \\
Refinement Effectiveness ($RE$) & --- & +9.5\% & +19.0\% & +19.0\% \\
\midrule
\textbf{Migration (Task B)} & & & & \\
Executed Successfully ($N=27$) & 23 (85.2\%) & 25 (92.6\%) & 26 (96.3\%) & 26 (96.3\%) \\
Passing on Clean SUT & 21 (77.8\%) & 24 (88.9\%) & 25 (92.6\%) & 25 (92.6\%) \\
Mean KSP (\%) & 90.7\% & 94.2\% & 96.8\% & 97.4\% \\
Cumulative Recovered Failures & --- & 2 / 4 (50.0\%) & 3 / 4 (75.0\%) & 3 / 4 (75.0\%) \\
Refinement Effectiveness ($RE$) & --- & +8.7\% & +13.0\% & +13.0\% \\
\bottomrule
\end{tabular}
\end{table}

Execution-guided refinement successfully resolved $66.7\%$ of initial failures in
Task A and $75.0\%$ of initial failures in Task B within two refinement iterations.
Most runtime failures were attributable to asynchronous event timing (e.g., typing
into the task input without triggering the \texttt{input} event required by client JS)
and race conditions during page reloads. Providing the execution error traceback
enabled the LLM to introduce explicit event dispatchers and auto-waiting assertions.

\subsection{RQ5: Baseline vs. Proposed Framework Comparison}
\label{subsec:rq5-results}

Table~\ref{tab:ablation-summary} summarizes the statistical comparison between
Baseline-1 and the Proposed Framework. Using the Wilcoxon signed-rank test across
paired test instances, the proposed framework demonstrates statistically significant
improvements on all dimensions ($p < 0.001$). The effect size calculated using
Cliff's delta ($\delta$) exceeds $0.70$ across all primary metrics, indicating
large, practically meaningful improvements.

\begin{table}[t]
\centering
\caption{Statistical Significance and Effect Sizes: Proposed Framework vs. Baseline-1.}
\label{tab:ablation-summary}
\small
\begin{tabular}{lcccc}
\toprule
\textbf{Metric} & \textbf{Baseline-1} & \textbf{Proposed} & \textbf{$p$-value (Wilcoxon)} & \textbf{Effect Size (Cliff's $\delta$)} \\
\midrule
Task A: ESR & 40.7\% & 92.6\% & $p < 0.001$ & 0.81 (Large) \\
Task A: Statement Coverage & 58.4\% & 94.8\% & $p < 0.001$ & 0.88 (Large) \\
Task B: ESR & 48.1\% & 96.3\% & $p < 0.001$ & 0.79 (Large) \\
Task B: KSP & 68.2\% & 97.4\% & $p < 0.001$ & 0.84 (Large) \\
Task B: SDR & 38.5\% & 3.7\% & $p < 0.001$ & -0.74 (Large) \\
\bottomrule
\end{tabular}
\end{table}

\subsection{Model-Level Comparison}
\label{subsec:model-comparison}

Evaluating performance across models reveals that Claude 3.5 Sonnet achieved the highest
first-pass migration success ($88.9\%$ initial $ESR$), followed closely by GPT-4o
($85.2\%$) and Gemini 1.5 Pro ($81.5\%$). However, after applying the validation-refinement
loop (Proposed), all three models converged to comparable final execution success
rates ($96.3\%$ for Claude 3.5 Sonnet and GPT-4o; $92.6\%$ for Gemini 1.5 Pro). This
demonstrates that the execution-guided refinement architecture effectively compensates
for initial model-specific code generation discrepancies.

% =========================================================
% 7. DISCUSSION
% =========================================================

\section{Discussion}
\label{sec:discussion}

The experimental findings offer critical insights into the dynamics of automated
test generation and cross-framework migration.

\subsection{The Peril of Silent Semantic Degradation}
A conventional evaluation that solely measures whether a migrated test compiles
and passes on the clean application would report an apparent pass rate of $76.9\%$
for direct prompting (Baseline-1). However, mutation analysis reveals that nearly
$40\%$ of those passing tests suffered from silent degradation: they dropped
assertions, checked trivial preconditions, or bypassed client-side validation logic,
thereby losing the ability to detect software regressions. By incorporating mutation
kill-set preservation ($KSP$) and silent degradation rate ($SDR$) into the evaluation
pipeline, our framework exposes and rectifies these critical semantic voids.

\subsection{Synchronization and Framework Idioms}
Our error analysis demonstrates that cross-framework UI test migration is fundamentally
not a simple syntax-mapping problem. Selenium tests frequently rely on synchronous
polling and page refreshes, whereas Playwright operates on a unified event loop with
automatic element actionability checks. Direct prompt-to-code translation consistently
mishandled input event listeners (\texttt{addEventListener('input')}), failing to trigger
client-side DOM counters unless explicit input sequences were dispatched. Supplying
framework-specific context ($C_F$) and structured action intents ($A_i$) resolved
these subtle behavioral mismatches.

\subsection{Iterative Refinement Convergence}
Refinement benefits plateau after two iterations ($k=2$). Across both tasks, initial
failures were largely corrected in the first refinement pass ($50$--$66\%$ recovery)
when the error traceback provided clear locator or timing diagnostics. Errors that
persisted beyond iteration 2 typically involved fundamental architectural misconceptions
(e.g., attempting to bypass server CSRF tokens without session establishment) rather
than minor syntax or selector errors.

\subsection{Threats to Validity}
\label{sec:threats}

\textbf{Construct Validity:} Construct validity concerns whether our metrics
accurately reflect test quality. We mitigated this by utilizing deterministic AST
and DOM mutation analysis, providing an objective, ground-truth measure of
fault-detection equivalence ($KSP$) rather than relying on superficial textual or
AST similarity.

\textbf{Internal Validity:} Variations in stochastic model outputs were mitigated
by repeating all experiments $K = 3$ times with a fixed low temperature ($0.2$) and
averaging results. Furthermore, our baseline check executed three identical pre-flight
runs to eliminate flaky tests prior to mutation evaluation.

\textbf{External Validity:} The evaluation was conducted on TaskBoard, an interactive
full-stack application. While TaskBoard encompasses authentic UI workflows (auth,
DOM updates, client-side validation, filtering), evaluating the framework across
larger enterprise codebases, additional language bindings (e.g., Java, TypeScript),
and other framework pairs (e.g., Cypress to Playwright) remains an important
direction for future work.

% =========================================================
% 8. CONCLUSION
% =========================================================

\section{Conclusion}
\label{sec:conclusion}

This paper presented a context-aware large language model framework for automated
test script generation and cross-framework migration. By structuring the transformation
pipeline through contextual acquisition, intermediate test-intent modeling ($Z$),
target script synthesis, and execution-guided iterative refinement, the framework
overcomes the critical limitations of conventional direct prompting.

We empirically evaluated the framework on an end-to-end benchmark comprising 27 test
scenarios and 57 deterministic mutants. The results show that the proposed framework
achieves $96.3\%$ execution success and $97.4\%$ kill-set preservation for
Selenium-to-Playwright migration, while dramatically reducing the silent degradation
rate from $38.5\%$ to $3.7\%$. For requirement-to-script generation, the framework
achieves $92.6\%$ execution success and $94.8\%$ statement coverage. These findings
provide compelling empirical evidence that execution feedback, contextual grounding,
and mutation-based semantic validation are essential for robust LLM-driven test automation.

% =========================================================
% REFERENCES
% =========================================================

\bibliography{sn-bibliography}

\end{document}
